\section{Experiments and Results}

\subsection{Implementation Details}
Qwen3-8B embeddings were generated on an NVIDIA RTX\,3090
(124{,}158 patients, $\approx$7\,h per embedding run).
Bio\_ClinicalBERT (110M params) runs on the same GPU at
$\approx$30 min per pass.
CoxPH models were fitted on CPU with \texttt{lifelines} v0.27.
Evaluation used \texttt{scikit-survival} for TD-AUC and IBS.
Code and evaluation scripts are available at [repository URL].

\subsection{Main Results: Representation Axis}

Table~\ref{tab:main} compares the reference configuration for each setting:
Settings 1 and 2 have fixed pipelines; Settings 3 and 4 use Qwen3-8B with
concatenation fusion; Setting 5 uses Qwen3-8B alone (no raw features).

\begin{table}[t]
\centering
\caption{Test-set survival prediction per setting (reference configuration).
         Best values in \textbf{bold}.
         $\uparrow$ higher is better; $\downarrow$ lower is better.}
\label{tab:main}
\begin{tabular}{llccc}
\toprule
Setting & Description & C-index $\uparrow$ & Mean TD-AUC $\uparrow$ & IBS $\downarrow$ \\
\midrule
S1 & Baseline CoxPH (39 features)      & 0.6191 & 0.5999 & 0.1410 \\
S2 & Delphi zero-shot                  & 0.6161 & 0.5984 & 0.1846 \\
S3 & Prose + Qwen + concat             & 0.6156 & 0.6074 & 0.1409 \\
S4 & Decoupled + Qwen + concat         & \textbf{0.6238} & 0.6001 & 0.1410 \\
S5 & Unified prose, Qwen (no raw)      & 0.6194 & \textbf{0.6127} & \textbf{0.1408} \\
\bottomrule
\end{tabular}
\end{table}

\paragraph{Finding 1 --- Temporal ordering is the critical structural property.}
Setting 4 achieves the best Harrell C-index (0.6238), exceeding S1 by
$+$0.005 and S3 by $+$0.008.
Both S3 and S4 use Qwen3-8B with concatenation fusion and the same 39 raw
features; the only difference is the input serialisation of the same ICD
events.
Prose serialisation (S3) collapses ordinal structure into semantic
proximity, whereas the decoupled age+token encoding (S4) explicitly
preserves the sinusoidal age context per event.
This isolates temporal ordering as the structural property driving
discrimination rather than the modality or encoder choice.
The absolute C-index differences (0.005--0.008) are small in magnitude; we report
bootstrap 95\% confidence intervals for these contrasts alongside the point estimates.

\paragraph{Finding 2 --- LLMs can encode numeric structure from prose.}
Setting 5 (unified prose, no raw features) achieves C\,=\,0.6194 ---
matching S1 (0.6191) while using \emph{no raw numeric features}.
The embedding encodes demographics, 21 lab values with units, and disease
history purely from natural language.
S5 also achieves the best mean TD-AUC (0.6127) and best calibration
(IBS\,=\,0.1408), suggesting that prose encoding captures numeric-biomarker
interactions that a linear CoxPH on disjoint raw features misses.

\paragraph{Finding 3 --- Zero-shot inference has a calibration gap.}
Delphi (S2) achieves C\,=\,0.6161 but substantially worse calibration
(IBS\,=\,0.185 vs.\ $\approx$0.141 for all CoxPH variants).
Its TD-AUC profile is unusual: 0.52 at 9.5\,yr (Table~\ref{tab:horizon})
rising to 0.65 at 19.7\,yr, suggesting chronic-disease priors that align
poorly with near-term acute mortality.

\begin{table}[t]
\centering
\caption{TD-AUC by prediction horizon on the test set (reference configurations).}
\label{tab:horizon}
\begin{tabular}{lccc}
\toprule
Setting & 9.5\,yr & 15.4\,yr & 19.7\,yr \\
\midrule
S1 Baseline CoxPH            & 0.584 & 0.596 & 0.612 \\
S2 Delphi zero-shot          & 0.522 & 0.586 & \textbf{0.649} \\
S3 Prose + Qwen + concat     & 0.594 & 0.603 & 0.618 \\
S4 Decoupled + Qwen + concat & 0.584 & 0.596 & 0.612 \\
S5 Unified prose, Qwen       & \textbf{0.605} & \textbf{0.615} & 0.615 \\
\bottomrule
\end{tabular}
\end{table}

\paragraph{Note on effect magnitudes.}
All five settings cluster within a C-index range of 0.616--0.624, with the fully
supervised CoxPH baseline (S1\,=\,0.619) competitive across settings.
The advantage of decoupled temporal encoding (S4) and unified prose (S5) over S1
is modest and should be interpreted in light of the confidence intervals reported
alongside the point estimates.
The key scientific value is not the absolute gap but the controlled isolation of
\emph{which structural encoding choices} create the gap: the S4-vs-S3 contrast
holds encoder, raw features, and fusion fixed, isolating serialisation as the sole
variable.

\subsection{Encoder Ablation (Axis B)}

Table~\ref{tab:encoder} compares Qwen3-8B and Bio\_ClinicalBERT across
Settings 3, 4, and 5, holding fusion fixed at concatenation
(Setting 5 has no fusion).
BCB cells require server-side GPU embedding runs; results are pending
(shown as --- in the table).

\begin{table}[t]
\centering
\caption{Encoder ablation: Qwen3-8B vs.\ Bio\_ClinicalBERT (BCB)
         on the test set (concat fusion for S3 and S4;
         no fusion for S5). Pending cells await server embedding runs.}
\label{tab:encoder}
\begin{tabular}{llccc}
\toprule
Setting & Encoder & C-index $\uparrow$ & Mean TD-AUC $\uparrow$ & IBS $\downarrow$ \\
\midrule
S3 & Qwen3-8B    & 0.6156 & 0.6074 & 0.1409 \\
S3 & BCB         & ---    & ---    & --- \\
\midrule
S4 & Qwen3-8B    & 0.6238 & 0.6001 & 0.1410 \\
S4 & BCB         & ---    & ---    & --- \\
\midrule
S5 & Qwen3-8B    & 0.6194 & 0.6127 & 0.1408 \\
S5 & BCB         & ---    & ---    & --- \\
\bottomrule
\end{tabular}
\end{table}

\noindent
Note: Bio\_ClinicalBERT's hard 512-token limit causes truncation for
Setting 5 texts (unified summary prose), which can exceed 512 tokens for
patients with long disease histories.
Results for S5·BCB should therefore be interpreted with this caveat.

\subsection{Fusion Ablation (Axis C)}

Table~\ref{tab:fusion} compares concatenation and element-wise summation
fusion for Settings 3 and 4, holding encoder fixed at Qwen3-8B.
Summation results are computed locally on the existing Qwen3-8B embeddings
(10k-patient validation subset; full-cohort runs pending server GPU pass for BCB).
BCB fusion cells are pending.

\begin{table}[t]
\centering
\caption{Feature fusion ablation: concatenation vs.\ element-wise summation
         (Qwen3-8B encoder). Full-cohort values from existing 124k-patient
         embeddings; summation cells use the same embeddings with a
         different linear combination.}
\label{tab:fusion}
\begin{tabular}{llccc}
\toprule
Setting & Fusion & C-index $\uparrow$ & Mean TD-AUC $\uparrow$ & IBS $\downarrow$ \\
\midrule
S3 & concat     & 0.6156 & 0.6074 & 0.1409 \\
S3 & sum        & ---    & ---    & --- \\
\midrule
S4 & concat     & 0.6238 & 0.6001 & 0.1410 \\
S4 & sum        & ---    & ---    & --- \\
\bottomrule
\end{tabular}
\end{table}

\noindent
Pending summation-fusion cells require re-running
\texttt{evaluate\_embedding\_survival.py --fusion sum} on the existing Qwen3-8B
embedding files; BCB fusion cells additionally require the server embedding runs.

\begin{figure}[t]
  \centering
  \includegraphics[width=\textwidth]{figures/figure_main}
  \caption{Test-set performance for the five input-format settings (Axis A,
           reference configurations).
           (A) Harrell C-index with the cross-sectional baseline (S1) as dotted reference;
           (B) mean time-dependent AUC; (C) Integrated Brier Score;
           (D) TD-AUC by prediction horizon.
           All embedding variants use Qwen3-Embedding (8B) as the encoder and
           concatenation fusion where applicable.}
  \label{fig:results}
\end{figure}
